\documentclass[11pt]{article}
\usepackage{mzqastyle}
\setrunningtitle{AI Analyst — Quantitative \& ML Models}

\begin{document}

\mzqatitlepage
  {Quantitative \& Machine-Learning Modeling Reference}
  {Quantitative \& Machine-Learning Models}
  {The complete mathematical specification}
  {This volume is the exhaustive mathematical reference for every quantitative and
   machine-learning model in AI Analyst. It specifies the cross-sectional
   gradient-boosted alpha model, the factor-structured risk models, the dual
   portfolio optimizers, the walk-forward out-of-sample backtest with
   Fama--French benchmarking, the rolling robust factor regression, the
   macro-cycle regime models (PCA/DFM, Gaussian HMM, temporal multimodal VAE, and
   regime-conditioned factor IC), and the deterministic valuation and scoring
   models --- each with its objective, estimator, algorithm, and the exact
   warehouse tables and code paths it runs on.

   \vspace{8pt}
   \textit{Companion volumes:} \emph{System Architecture \& Operations},
   \emph{qlib Integration}, and \emph{SEC \& EDINET Metrics Mapping \&
   Reconciliation}.}
  {Modeling Reference \quad·\quad Generated 31 July 2026 \quad·\quad Independent research, not investment advice.}

\mzqatoc
\clearpage

%======================================================================
\part*{Part I\quad Foundations}
\addcontentsline{toc}{section}{Part I\quad Foundations}

\section{Scope and notation}

AI Analyst answers two different questions with two different families of model.
The \emph{valuation} models (Part IV) answer \emph{``what is this business worth?''}
deterministically, from projected cash flows and a cost of capital. The
\emph{quantitative} models (Parts II--III) answer \emph{``what should I expect from
this stock, and how should I size it?''} statistically, from a learned
cross-sectional return forecast, a forward covariance, and a portfolio optimizer,
conditioned on an estimated macro regime. This document is the mathematical
specification of both families.

\begin{specbox}
\tabhead{Notation.}\; $\wv\in\R^N$ portfolio weights; $\muv\in\R^N$ expected
(annualized) returns; $\Sigma\in\R^{N\times N}$ return covariance; $r_f$ the
risk-free rate; $\ones$ the ones vector; $\odot$ the elementwise (Hadamard)
product; $z$ a cross-sectional $z$-score. Time subscripts $t$ index months unless
stated otherwise. Annualization uses $252$ trading days for daily series and $12$
for monthly series. A ``\,hat\,'' ($\hat{\,\cdot\,}$) denotes an estimate or a
model prediction.
\end{specbox}

Every model is designed to \emph{degrade gracefully}: if a trained artifact,
an optional dependency (\code{hmmlearn}, PyTorch, a MIP solver), or a warehouse
table is absent, the code returns a defined fallback rather than failing, so the
application always runs. Those fallbacks are stated with each model.

\section{The data substrate: a point-in-time fundamentals panel}
\label{sec:panel}

Every cross-sectional model consumes the same object: a monthly panel of
standardized fundamentals aligned \emph{point-in-time} (PIT) so that no signal
ever uses data that was not yet public.

\subsection{Point-in-time alignment}
A fundamental value with fiscal-period end $p$ is stamped usable only at
\begin{equation}
  \tau(p) \;=\; \operatorname{monthend}\!\big(p + 90\text{ days}\big),
\end{equation}
a conservative proxy for the filing lag. In SQL the alignment is
\code{(date\_trunc('month', period\_end + INTERVAL '90 days') + INTERVAL '1 month - 1 day')}.
The alpha model (\S\ref{sec:alpha}) and the regime-IC engine (\S\ref{sec:regic})
share this one loader, so the research signal and the production forecast see an
identical, lookahead-safe history.

\subsection{Feature families and cross-sectional standardization}
\label{sec:standardize}
Features are importance-ranked standardized fundamentals
(\code{fact\_metrics\_\{us,jp\}}, keeping \code{importance}\,$\le 2$) drawn from
four families:

\begin{center}\small
\begin{tabular}{@{}>{\bfseries\color{navy}}l p{10.8cm}@{}}
\toprule
\tabhead{Family} & \tabhead{Representative factors (substring patterns)} \\
\midrule
value           & book-to-market, earnings yield, FCF yield, dividend, EV/\,$\cdot$, valuation multiples \\
quality         & return-on-\,$\cdot$, ROE/ROA, margins, accruals, cash conversion, asset turnover \\
growth          & revenue/earnings growth, investment, capex, R\&D, sales change \\
market\_factor  & market beta, volatility, momentum, residual/abnormal return \\
\bottomrule
\end{tabular}
\end{center}

For each factor $j$, name $i$, and month $t$, the raw value $x^{(j)}_{i,t}$ is
\emph{cross-sectionally} standardized and winsorized (mirroring qlib's
\code{CSZScoreNorm}):
\begin{equation}
  z^{(j)}_{i,t} \;=\; \clip\!\left(
    \frac{x^{(j)}_{i,t}-\bar\mu^{(j)}_t}{\hat\sigma^{(j)}_t},\,-3,\,3\right),
  \qquad
  \bar\mu^{(j)}_t=\frac1{N_t}\sum_{i} x^{(j)}_{i,t},\quad
  \hat\sigma^{(j)}_t=\Std_i\, x^{(j)}_{i,t}.
\end{equation}
Standardizing within each date makes factors comparable across names and dates
and robust to per-metric scale; a zero-variance cross-section maps to $0$
(the neutral value). A missing feature is imputed to the cross-sectional mean,
i.e.\ $z=0$.

\subsection{Label}
The supervised label is the forward one-month (or three-month) total return,
\begin{equation}
  y_{i,t} \;=\; r_{i,\,t\to t+h}, \qquad h\in\{1,3\}\ \text{months}.
\end{equation}
The panel is built as a qlib MultiIndex frame --- rows $(\text{datetime},
\text{instrument})$, columns $(\text{feature}\,|\,\text{label})$ --- consumed
directly by \code{DataHandlerLP.from\_df}. JP fundamentals carry a \code{.T}
suffix that is stripped so fundamentals join prices. Thin cross-sections
($<30$ names) are dropped as unstable for ranking and normalization.

\begin{keybox}[Design]
The panel is the single contract between the warehouse and the models: a monthly,
cross-sectional, lookahead-safe frame of $z$-scored fundamentals plus a forward
return. Expected returns from the alpha model are therefore \emph{monthly}
forecasts --- annualize by $\times\,12$ before feeding a portfolio optimizer.
\end{keybox}

%======================================================================
\part*{Part II\quad The qlib quant desk}
\addcontentsline{toc}{section}{Part II\quad The qlib quant desk}

\section{Cross-sectional alpha model (return prediction)}
\label{sec:alpha}
\emph{Code:} \code{api/quant/qlib\_alpha.py}, \code{api/quant/qlib\_data.py},
\code{api/quant/alpha\_signal.py}; built on qlib's \code{LGBModel} (LightGBM).

\subsection{Objective}
Learn a function $f:\zv\mapsto\hat y$ mapping the standardized cross-sectional
feature vector of a name to its expected forward return. The prediction
$\hat y=f(\zv)$ is the \emph{alpha signal}: a per-name, per-month score
interpreted as an expected forward return.

\subsection{Gradient-boosted decision trees}
The model is an additive ensemble of $M$ regression trees with learning rate
$\eta$,
\begin{equation}
  F_M(\zv)\;=\;\sum_{m=1}^{M}\eta\,h_m(\zv),
  \qquad h_m\ \text{a regression tree.}
\end{equation}
Under the squared-error objective
$\mathcal L=\sum_i\tfrac12\big(y_i-F(\zv_i)\big)^2$, boosting stage $m$ fits the
new tree $h_m$ to the negative functional gradient --- here simply the residual ---
with unit Hessian:
\begin{equation}
  g_i \;=\; -\frac{\partial \mathcal L}{\partial F}\bigg|_{F_{m-1}}
        \;=\; y_i-F_{m-1}(\zv_i),
  \qquad h_i \;=\; \frac{\partial^2 \mathcal L}{\partial F^2}\;=\;1.
\end{equation}
LightGBM grows each tree \emph{leaf-wise} over a histogram approximation of the
features. Splitting a node with instance set $I$ into $I_L,I_R$ yields the gain
\begin{equation}
  \text{Gain}=\tfrac12\!\left[
  \frac{\big(\sum_{i\in I_L}g_i\big)^2}{\sum_{i\in I_L}h_i+\lambda}
  +\frac{\big(\sum_{i\in I_R}g_i\big)^2}{\sum_{i\in I_R}h_i+\lambda}
  -\frac{\big(\sum_{i\in I}g_i\big)^2}{\sum_{i\in I}h_i+\lambda}\right]-\gamma,
\end{equation}
where $\lambda$ is the $L_2$ leaf regularizer and $\gamma$ the per-split penalty.
Trees capture non-linearities and factor \emph{interactions} that a linear factor
model cannot represent.

\begin{modelbox}{Shipped hyperparameters (\texttt{DEFAULT\_LGB\_PARAMS})}
\small
\begin{tabular}{@{}llll@{}}
\code{loss} = mse & \code{learning\_rate} = 0.03 & \code{num\_leaves} = 31 & \code{max\_depth} = 6\\
\code{colsample\_bytree} = 0.8 & \code{subsample} = 0.8 & \code{subsample\_freq} = 1 & \code{min\_child\_samples} = 50\\
\code{lambda\_l1} = 1.0 & \code{lambda\_l2} = 1.0 & \code{num\_boost\_round} = 500 & \code{early\_stopping\_rounds} = 40\\
\end{tabular}\\[4pt]
Regularization ($L_1/L_2$ leaf penalties, shallow depth, feature/row subsampling)
plus early stopping on a held-out validation segment control overfitting in a
setting with a low signal-to-noise ratio.
\end{modelbox}

\subsection{Training, splitting, and persistence}
The panel is split \emph{chronologically} by unique dates into contiguous
train / validation / test segments (defaults $70/15/15$). Training runs the real
qlib \code{DatasetH} pipeline; the \code{AlphaLGB} subclass overrides \code{fit}
to attach LightGBM's early-stopping and evaluation callbacks directly, so training
runs with plain LightGBM and needs no \code{qlib.init} or metric-tracking backend.
Prediction runs the stored booster on any reindexed feature matrix --- including a
single live month --- so scoring never rebuilds a \code{DatasetH}. Artifacts (the
booster plus metadata: feature identity/order, label horizon, train range, IC)
are serialized per jurisdiction and hot-reloaded on file mtime.

For the live cross-section the model builds a \emph{features-only} panel
(\code{require\_label=False}) over a short lookback and scores the newest month:
the current month has no realized forward return, so a label join would drop
exactly the cross-section we want to score.

The split is contiguous by default, which leaks whenever the label horizon exceeds one
month: an $h$-month label observed at the last training month is not realized until
$h$ months later, i.e.\ inside the validation block. \code{time\_segments} therefore
accepts an \code{embargo\_months} gap that purges that many months from the end of each
earlier segment; the research loop of \S\ref{sec:research} sets it to the label horizon
by default. Figures computed without it are optimistic and are not comparable with the
purged walk-forward of \S\ref{sec:purged}.

\subsection{Evaluation: the information coefficient}
Predictive power is measured per month by the cross-sectional correlation between
prediction and realized return:
\begin{align}
  \IC_t &= \Corr\big(\hat y_{\cdot,t},\,y_{\cdot,t}\big)\ \text{(Pearson)}, &
  \RankIC_t &= \rho_{\text{Spearman}}\big(\hat y_{\cdot,t},\,y_{\cdot,t}\big),
\end{align}
summarized by the mean IC and the information ratio of the signal,
$\mathrm{ICIR}=\overline{\IC}/\Std(\IC)$ (and its rank analogue). A monthly equity
signal with $\overline{\RankIC}\approx 0.05\text{--}0.10$ is usable. The
\emph{in-sample} test segment of the shipped US/JP models reaches
$\overline{\RankIC}\approx 0.08\text{--}0.14$; the \emph{out-of-sample}
walk-forward figure (\S\ref{sec:bt}) is materially lower and is the number used
for any performance claim.

Note that $\mathrm{ICIR}$ as defined above is a bare mean/standard-deviation ratio, not
the annualized information ratio the name usually denotes; the comparable monthly-panel
figure carries a $\sqrt{12}$ scaling (\S\ref{sec:battery}). Both are reported by
\code{evaluate()}.

\section{Factor-structured risk models (forward covariance)}
\label{sec:risk}
\emph{Code:} \code{api/quant/qlib\_risk.py} (statistical) and
\code{api/quant/risk.py} (Ledoit--Wolf, Fama--French).

A portfolio optimizer needs a forward covariance $\Sigma$ that is well conditioned
and invertible. A raw $N\times N$ sample covariance from a short daily window is
noisy and near-singular. All three risk models replace it with a structured,
lower-parameter estimator. Each is selectable at runtime.

\subsection{Statistical factor covariance (qlib \texttt{StructuredCovEstimator})}
Returns are modeled as driven by $K$ latent statistical factors,
\begin{equation}
  X = B\,F\tr + U, \qquad
  \Cov(X) = F\,\Cov(B)\,F\tr + \diag\!\big(\Var(U)\big),
\end{equation}
where $X\in\R^{T\times N}$ is the return panel, $F\in\R^{N\times K}$ the factor
\emph{exposures} (principal components; factor analysis is also supported),
$B\in\R^{T\times K}$ the factor \emph{returns} (component scores), and $U$ the
idiosyncratic residual. Writing $\Sigma_B=\Cov(B)$ and $\mathbf d=\Var(U)$,
\begin{equation}
  \boxed{\;\Sigma \;=\; F\,\Sigma_B\,F\tr + \diag(\mathbf d)\;}
\end{equation}
--- a low-rank-plus-diagonal shrinkage of the sample covariance. The number of
factors is clamped to $K=\min(K_0,\,N-1,\,T-2)$ for safety. Estimation is on
\emph{raw} returns (the estimator's percentage-rescaling is disabled, keeping
variances in return units); $\Sigma_B$ and $\mathbf d$ are annualized by $\times
252$, and $\Sigma$ is rebuilt from the annualized decomposition, symmetrized, and
given a $10^{-10}$ ridge for numerical positive-definiteness. Forward volatility
per name is $\sigma_i=\sqrt{\Sigma_{ii}}$. The decomposition $(F,\Sigma_B,\mathbf
d)$ is consumed directly by the enhanced-indexing optimizer (\S\ref{sec:enh}).

\subsection{Ledoit--Wolf shrinkage (constant-correlation target)}
\label{sec:lw}
An analytic shrinkage of the MLE sample covariance $S$ toward a
constant-correlation target $F$. With variances $s_{ii}$, standard deviations
$\sigma_i=\sqrt{s_{ii}}$, and mean pairwise correlation
$\bar r=\tfrac{2}{N(N-1)}\sum_{i<j}s_{ij}/(\sigma_i\sigma_j)$, the target is
$F_{ii}=s_{ii}$, $F_{ij}=\bar r\,\sigma_i\sigma_j$, and
\begin{equation}
  \Sigma^\star=\delta F+(1-\delta)S, \qquad
  \delta=\clip\!\Big(\tfrac{\pi/\gamma}{T},\,0,\,1\Big),
\end{equation}
where $\gamma=\lVert S-F\rVert_F^2$ measures target misspecification and
$\pi=\sum_{i,j}\big[\tfrac1T\sum_t (x_{ti}x_{tj}-s_{ij})^2\big]$ is the summed
asymptotic variance of the entries of $S$. The optimal intensity $\delta$ trades
sample noise against target bias; the result is annualized $\times 252$.

\subsection{Fama--French fundamental covariance}
\label{sec:ffcov}
A named-factor analogue of the statistical model:
\begin{equation}
  \Sigma = B\,F_{\text{ann}}\,B\tr + D,
\end{equation}
with $B$ the per-name FF loadings (\code{fact\_factor\_loadings}, \S\ref{sec:ffreg}),
$F_{\text{ann}}=252\cdot\Cov(\text{factor returns})$ estimated over the lookback,
and $D=\diag(\text{residual\_vol}_i^2)$ from the regression meta. Model ids
\code{ff3}, \code{ff5}, \code{ff5\_plus\_mom} select the factor set
(FF6 $=$ FF5 $+$ momentum). Names with missing loadings fall back to their sample
diagonal with a surfaced warning; an unavailable factor model degrades to
Ledoit--Wolf.

\section{Portfolio optimization (dual backend)}
\label{sec:opt}
\emph{Code:} \code{api/quant/qlib\_optimize.py} dispatches at runtime to a
\emph{native} optimizer (\code{api/quant/optimize.py}) or any qlib optimizer;
both consume the same $\muv,\Sigma$ and return a uniform \code{PortfolioSolution}.

\subsection{Native optimizer: multi-term mean--variance utility}
The native backend maximizes a mean--variance utility with factor-tilt,
turnover, and concentration penalties:
\begin{equation}
  \max_{\wv}\;\;
  \wv\tr\muv
  -\tfrac12\lambda\,\wv\tr\Sigma\wv
  -\gamma_f\big\lVert S\odot a\odot(B\tr\wv-\mathbf t)\big\rVert_2^2
  -\gamma_{\text{to}}\lVert\wv-\wv_0\rVert_1
  -\gamma_c\lVert\wv\rVert_2^2,
\end{equation}
subject to full investment $\ones\tr\wv=1$, per-name bounds $l_i\le w_i\le u_i$,
a gross cap $\sum_i|w_i|\le G$, a short-leg cap $\sum_i\max(-w_i,0)\le \rho\sum_i|w_i|$,
sector caps $\mathbf a_s\tr\wv\le c_s$, an annualized volatility cap
$\wv\tr\Sigma\wv\le\sigma_{\max}^2$, and factor caps/ranges
$|\mathbf b_j\tr\wv|\le\text{cap}_j$. Here $\lambda$ is risk aversion; $\gamma_f$
penalizes deviation of the portfolio's factor exposures $B\tr\wv$ from targets
$\mathbf t$ (scaled by $S$, gated by the active mask $a$); $\gamma_{\text{to}}$
penalizes turnover against the current book $\wv_0$; and $\gamma_c$ penalizes
concentration.

The problem is solved by sequential least-squares programming (SLSQP), which
handles the nonlinear volatility and factor constraints natively and converges
quickly at portfolio scale ($N\lesssim100$). SLSQP is given the analytic
gradient
\begin{equation}
  \nabla \;=\; -\muv+\lambda\Sigma\wv
  +2\gamma_f\,B\big(S^2\odot a\odot(B\tr\wv-\mathbf t)\big)
  +\gamma_{\text{to}}\,\sgn(\wv-\wv_0)+2\gamma_c\wv .
\end{equation}

\paragraph{Cardinality (mixed-integer program).} When a maximum position count
$\lVert\wv\rVert_0\le K$ is requested, the problem becomes a mixed-integer program
with binary inclusion indicators $y_i\in\{0,1\}$:
\begin{equation}
  \sum_i y_i\le K,\qquad
  \begin{cases}
    0\le w_i\le u_i\,y_i, & \text{long-only,}\\[2pt]
    -u_i\,y_i\le w_i\le u_i\,y_i, & \text{long/short,}
  \end{cases}
  \qquad w_i\ge f\,y_i\ \text{(floor on included names),}
\end{equation}
retaining the convex volatility, sector, and factor constraints. It is solved via
cvxpy over the first available MIP backend (HiGHS $\to$ SCIP $\to$ CBC $\to$
GLPK\_MI $\to$ MOSEK); if none is installed the route surfaces a clear
``MIP unavailable'' error rather than silently dropping the constraint.

\paragraph{Diagnostics and frontier.} The solver reports factor exposures
$B\tr\wv$, marginal risk contributions $(\Sigma\wv)\odot\wv/\sqrt{\wv\tr\Sigma\wv}$,
the effective number of names $1/\sum_i w_i^2$, and the covariance condition
number; it also sweeps $\lambda$ over a fixed grid to trace a mini efficient
frontier.

\subsection{qlib optimizers (long-only, fully invested)}
All satisfy $\ones\tr\wv=1,\ \wv\ge0$:
\begin{center}\small
\begin{tabular}{@{}>{\bfseries\color{navy}}l l@{}}
\toprule
\tabhead{Backend} & \tabhead{Objective} \\
\midrule
GMV (min variance)   & $\displaystyle\min_{\wv}\ \wv\tr\Sigma\wv$ \\[3pt]
MVO (mean--variance) & $\displaystyle\max_{\wv}\ \lambda\,\tilde r\tr\wv-\wv\tr\Sigma\wv$
   \quad ($\tilde r$ scaled to $\Sigma$'s volatility) \\[3pt]
RP (risk parity)     & choose $\wv$ s.t.\ $w_i(\Sigma\wv)_i=w_j(\Sigma\wv)_j\ \forall i,j$ \\[3pt]
Inverse-vol          & $\displaystyle w_i=\frac{1/\sigma_i}{\sum_j 1/\sigma_j},\quad \sigma_i=\sqrt{\Sigma_{ii}}$ \\
\bottomrule
\end{tabular}
\end{center}
Risk parity equalizes the risk contributions $\mathrm{RC}_i=w_i(\Sigma\wv)_i$,
which sum to total variance $\sum_i\mathrm{RC}_i=\wv\tr\Sigma\wv$; GMV/MVO/RP are
solved with \code{scipy.optimize} and admit optional turnover ($\delta$) and $L_2$
($\alpha$) terms.

\subsection{Enhanced indexing (benchmark-relative)}
\label{sec:enh}
Given benchmark weights $\wv_b$, active weights $\mathbf d=\wv-\wv_b$, and active
factor exposure $\mathbf v=\mathbf d\tr F$, the enhanced-indexing optimizer
maximizes active return minus active (tracking-error) risk under bounded name and
factor deviations:
\begin{equation}
  \max_{\wv}\;\;
  \mathbf d\tr\mathbf r
  -\lambda\big(\mathbf v\tr\Sigma_B\mathbf v+\mathbf d\tr\diag(\mathbf d_u)\,\mathbf d\big)
  \quad\text{s.t.}\quad
  \wv\ge0,\ \ones\tr\wv=1,\
  \lVert\wv-\wv_0\rVert_1\le\delta,\
  |\mathbf d|\le b_{\text{dev}},\ |\mathbf v|\le f_{\text{dev}},
\end{equation}
where $\mathbf d_u=\Var(U)$ is the specific variance. It is a convex quadratic
program consuming the statistical-covariance decomposition $(F,\Sigma_B,\mathbf
d_u)$ of \S\ref{sec:risk} directly, solved by cvxpy/ECOS; on solver failure it
degrades to the benchmark weights with a warning.

\subsection{Reported portfolio statistics}
For a solution $\wv^\star$ the desk reports
$\E[r]_{\text{ann}}=\wv^{\star\mathsf T}\muv$,
$\sigma_{\text{ann}}=\sqrt{\wv^{\star\mathsf T}\Sigma\wv^\star}$, Sharpe
$=(\E[r]_{\text{ann}}-r_f)/\sigma_{\text{ann}}$, and the factor exposures and
risk contributions above.

\section{Walk-forward out-of-sample backtest}
\label{sec:bt}
\emph{Code:} \code{api/quant/qlib\_backtest.py}. A cross-sectional signal
backtest that is out-of-sample by construction, benchmarked against Fama--French.

\subsection{Walk-forward predictions}
Let $t_0<t_1<\dots<t_T$ be the panel months, restricted to an investable universe
(market cap $\ge\$2$B) with each month's forward returns winsorized to the
$1$st/$99$th cross-sectional percentile (removing micro-cap data artifacts). For
each test month $t_i$ with $i\ge m_{\min}$ (default $m_{\min}=12$; refit every
$q=3$ months) a booster $g_{t_i}$ is trained on \emph{all prior} months
$\{t_j:t_j<t_i\}$ and predicts $\hat\alpha_{k,t_i}=g_{t_i}(\zv_{k,t_i})$. No month
is ever scored by a model that trained on it. The expensive walk-forward
predictions are cached separately from portfolio construction, so re-tuning
$\text{top-}k$ or long/short re-ranks instantly.

\subsection{Portfolio construction}
Ranking names by $\hat\alpha$ each month, the equal-weight books realize
\begin{equation}
  r^{\text{L}}_{t}=\frac1k\!\!\sum_{i\in\text{top-}k}\!\! r_{i,t},
  \qquad
  r^{\text{LS}}_{t}=\frac1k\!\!\sum_{i\in\text{top-}k}\!\! r_{i,t}
  -\frac1k\!\!\sum_{i\in\text{bot-}k}\!\! r_{i,t},
\end{equation}
for the long-only and long/short strategies respectively; a month is skipped if
it has fewer than $\max(k,10)$ (long) or $2k$ (long/short) investable names.

\subsection{Performance and risk metrics}
With monthly returns $\{r_t\}_{t=1}^n$, equity $E_t=\prod_{s\le t}(1+r_s)$, and
monthly risk-free $r_{f,t}$:
\begin{align}
  R_{\text{ann}}&=\Big(\textstyle\prod_t(1+r_t)\Big)^{12/n}-1, &
  \sigma_{\text{ann}}&=\Std(r_t)\sqrt{12}, \\
  \text{Sharpe}&=\frac{R_{\text{ann}}-R^{f}_{\text{ann}}}{\sigma_{\text{ann}}}, &
  \text{Sortino}&=\frac{R_{\text{ann}}-R^{f}_{\text{ann}}}{\Std(r_t:\,r_t<0)\sqrt{12}}, \\
  \text{MaxDD}&=\min_t\Big(\tfrac{E_t}{\max_{s\le t}E_s}-1\Big), &
  \text{Hit}&=\tfrac1n\textstyle\sum_t \mathbf 1[r_t>0].
\end{align}

\subsection{Fama--French benchmarking and factor attribution}
The benchmark is the FF \emph{market} total return
$r^m_t=(\text{Mkt-RF})_t+(\text{RF})_t$, compounded from daily
\code{fact\_fama\_french} to monthly and aligned to $\{r_t\}$. Relative to it,
\begin{equation}
  \text{excess}=R_{\text{ann}}-R^m_{\text{ann}},\quad
  \text{TE}=\Std(r_t-r^m_t)\sqrt{12},\quad
  \text{IR}=\frac{\text{excess}}{\text{TE}},\quad
  \beta=\frac{\Cov(r,r^m)}{\Var(r^m)} .
\end{equation}
To separate genuine alpha from factor exposure, the strategy's excess return is
regressed on the FF five factors plus momentum,
\begin{equation}
  r_t-r_{f,t}=\alpha
  +\beta_{\text{Mkt}}f^{\text{Mkt}}_t+\beta_{\text{SMB}}f^{\text{SMB}}_t
  +\beta_{\text{HML}}f^{\text{HML}}_t+\beta_{\text{RMW}}f^{\text{RMW}}_t
  +\beta_{\text{CMA}}f^{\text{CMA}}_t+\beta_{\text{Mom}}f^{\text{Mom}}_t+\varepsilon_t .
\end{equation}
OLS on the design $A=[\ones\ \ X]$ gives $\hat\theta=(A\tr A)^{-1}A\tr y$, residual
variance $\hat s^2=\lVert\varepsilon\rVert^2/(n-K-1)$, coefficient covariance
$\hat s^2(A\tr A)^{-1}$, the alpha $t$-statistic
$t_\alpha=\hat\alpha/\mathrm{se}(\hat\alpha)$, annualized alpha
$\alpha_{\text{ann}}=(1+\hat\alpha)^{12}-1$, and
$R^2=1-\lVert\varepsilon\rVert^2/\sum_t(y_t-\bar y)^2$.

\begin{keybox}[Interpretation]
The walk-forward design is what makes the backtest honest: because every month is
scored by a model that never saw it, the out-of-sample $\overline{\RankIC}$ and
$t_\alpha$ are the figures used for any claim about the signal. They are markedly
weaker than the in-sample test-segment numbers of \S\ref{sec:alpha} --- the
expected outcome of removing look-ahead, and the reason performance is always
quoted from this path.
\end{keybox}

%======================================================================
\part*{Part III\quad Factor and macro-regime models}
\addcontentsline{toc}{section}{Part III\quad Factor and macro-regime models}

%======================================================================
\section{Agentic model research}
\label{sec:research}

\emph{Code:} \code{api/quant/research/}. The estimators in this section replace the
single fixed specification of \S\ref{sec:alpha} with a searched one. The search is
driven by four LLM agents (their orchestration is documented in the \emph{qlib
Integration} volume); what follows is the \emph{statistics} they steer.

\subsection{The searched specification}
A \code{TrainingSpec} is a tuple $\theta$ collecting every choice that was previously
hard-coded: the sample filters, the outlier treatment, the standardization estimator,
the imputation rule, the feature block, the split geometry, the model family and its
hyperparameters. Writing $\mathcal D$ for the raw panel, a round evaluates
\begin{equation}
  \theta \;\longmapsto\; \big(\hat g_\theta,\ \mathcal M(\theta)\big),
  \qquad
  \hat g_\theta = \text{fit}\big(\Pi_\theta(\mathcal D)\big),
\end{equation}
where $\Pi_\theta$ is the preprocessing map below and $\mathcal M$ the metric battery
of \S\ref{sec:battery}. Agents propose a \emph{patch} to $\theta$; a deterministic
validator whitelists every key and clamps every value before it can reach a fit, so
the reachable set is exactly the specified product space.

\subsection{Robust cross-sectional estimators}
The incumbent standardization (\S\ref{sec:standardize}) uses the cross-sectional mean
and standard deviation, then clips at $\pm3$. Both moments are computed \emph{from the
sample that contains the outlier}, so a single extreme value shifts the centre and
inflates the scale for every other name in that month; the clip then protects the
model from the outlier but not the cross-section from its influence. Three
alternatives are searchable.

\paragraph{Winsorization by order statistics.}
For a month with $N_t$ observations and a tail fraction $q$, let $k=\max(1,\lfloor
qN_t\rfloor)$ and let $x_{(1)}\le\dots\le x_{(N_t)}$ be the order statistics. Each of
the $k$ most extreme observations per tail is replaced by the next one in:
\begin{equation}
  \tilde x_{i,t} \;=\; \clip\big(x_{i,t},\; x_{(k+1)},\; x_{(N_t-k)}\big).
\end{equation}
This is deliberately \emph{not} $\clip(x,\hat Q_q,\hat Q_{1-q})$ with interpolated
sample quantiles. An interpolated quantile is estimated from the contaminated sample,
so on a thin cross-section the outlier drags its own bound up to meet it --- with
$N_t=60$ and $q=0.01$ a planted value of $500$ pulled $\hat Q_{0.99}$ to $\approx206$
and the clip did essentially nothing. The order-statistic form has no such feedback,
and the floor $k\ge1$ keeps it meaningful when $qN_t<1$.

\paragraph{Robust standardization.}
With $\mathrm{med}_t$ the cross-sectional median and
$\mathrm{MAD}_t=\mathrm{med}_i\,|x_{i,t}-\mathrm{med}_t|$,
\begin{equation}
  z_{i,t} \;=\; \clip\!\left(\frac{x_{i,t}-\mathrm{med}_t}{1.4826\,\mathrm{MAD}_t},
  \,-c,\,c\right),
\end{equation}
the constant $1.4826$ making the MAD a consistent estimator of $\sigma$ under
normality. Both the location and the scale have a $50\%$ breakdown point, against
$0\%$ for the mean/standard-deviation pair.

\paragraph{Rank transform.}
Discarding magnitude entirely and keeping only the ordering the model is scored on:
\begin{equation}
  z_{i,t} \;=\; 2\left(\frac{\rank_i(x_{i,t})}{N_t}-\tfrac12\right)\in[-1,1].
\end{equation}

\paragraph{Imputation and neutralization.}
A missing feature may be filled at the cross-sectional \emph{median} rather than the
mean (the incumbent's $z=0$), or the row dropped. Features may additionally be
residualized within each month against sector dummies and $\log$ size,
\begin{equation}
  \tilde{\mathbf z}_t \;=\; \mathbf z_t - Z_t\big(Z_t\tr Z_t\big)^{-1}Z_t\tr\mathbf z_t ,
\end{equation}
one least-squares solve per month for the whole feature block, which asks whether a
factor says anything beyond which industry a company is in.

\subsection{Purged walk-forward evaluation}
\label{sec:purged}
The incumbent split is contiguous with no gap, so with an $h$-month label the last $h$
training months share realized future returns with the validation block. Let $E$ be
the embargo, defaulting to $E=h$. Training at prediction month $t_i$ is restricted to
\begin{equation}
  \mathcal T_{t_i} \;=\; \{\,t_j \;:\; t_j \le t_i - E\,\},
\end{equation}
so a model predicting month $t$ has seen only labels already realized by $t-E$ --- the
constraint a live deployment actually operates under. The same gap purges the
train/validation seam of the inner split used for early stopping.

\begin{notebox}[The embargo is not cosmetic]
On the shipped US \code{forward\_12m} model the purged walk-forward rank-IC of the
incumbent specification is \emph{negative}. Figures quoted from the unpurged
single-block split are not comparable with it and are optimistic by construction.
\end{notebox}

\subsection{The metric battery}
\label{sec:battery}
Organized by quality attribute rather than by convenience, after Lewis et al.'s
observation that ML evaluation overwhelmingly reports predictive accuracy and little
else, and that testing across attributes surfaces a wider class of defects.

\paragraph{Out-of-sample $R^2$.}
Both conventions are reported, because they diverge and quoting one without saying
which is how implausible figures arise:
\begin{equation}
  R^2_{\text{oos},0} = 1-\frac{\sum (y-\hat y)^2}{\sum y^2},
  \qquad
  R^2_{\text{oos},\mu} = 1-\frac{\sum (y-\hat y)^2}{\sum (y-\bar y)^2}.
\end{equation}
The zero-benchmarked form is the return-prediction standard: the historical mean is
itself a fitted quantity, and benchmarking against it flatters the model.

\paragraph{Scaled information ratio.}
$\mathrm{ICIR}=\overline{\IC}/\Std(\IC)$ is a bare ratio, not the annualized
information ratio the name suggests. On a monthly panel the comparable figure is
\begin{equation}
  \mathrm{ICIR}^{\text{ann}} \;=\; \frac{\overline{\IC}}{\Std(\IC)}\sqrt{12},
\end{equation}
roughly $3.46\times$ larger. Both are reported; the unscaled key retains its meaning
for existing consumers.

\paragraph{Uncertainty.}
A percentile bootstrap over the per-month $\RankIC_t$ series gives a $95\%$ interval
for $\overline{\RankIC}$. An interval spanning zero is treated as \emph{no
demonstrated skill} regardless of the point estimate, and blocks promotion.

\paragraph{Ranking quality.}
Beyond the top-minus-bottom decile spread and its $t$-statistic, the \emph{monotonicity}
$\rho_{\text{Spearman}}(d,\bar r_d)$ between decile index $d$ and its mean realized
return distinguishes a genuine cross-sectional ordering from a tail effect in a
handful of names --- a large spread with low monotonicity is the latter.

\paragraph{Implementability and hygiene.}
Turnover is the fraction of the top-$k$ book replaced each month. The FF5$+$momentum
regression of the long/short series separates alpha from factor exposure; a small
$t_\alpha$ beside a large $R^2$ means the strategy is a factor tilt. The
\emph{factor-neutral} rank-IC re-scores the signal after residualizing it
cross-sectionally on the names' Fama--French betas.

\paragraph{Stability.}
The train-minus-OOS IC gap is the overfitting tell. Feature-attribution rankings are
themselves checked for stability across refits (top-$k$ Jaccard and rank correlation)
before any conclusion is drawn from them, and the dispersion of the selected model
complexity across windows is reported: hyperparameters that swing between windows
indicate a low signal-to-noise ratio rather than a tuning opportunity.

\subsection{Perturbation robustness rating}
Adapted from the causal rating method of Lakkaraju et al.\ for time-series
forecasters. The fitted model is held \emph{fixed} and the data it is served is
degraded, so the procedure needs no access to the training pipeline and is auditable
by a third party. Five perturbations map to concrete warehouse failure modes: feature
values dropped (missing or late filings), halved (unit-scale and concept-mapping
errors), shifted one month stale (point-in-time slippage), the newest month withheld
(reporting delay), and fat-tail contamination of the label (corporate actions).

For perturbation $P$ the relative degradation is
\begin{equation}
  \delta_P \;=\; \frac{\overline{\RankIC}_{\varnothing}-\overline{\RankIC}_{P}}
                      {\big|\overline{\RankIC}_{\varnothing}\big|}.
\end{equation}
GICS industry is the confounder: a perturbation may appear damaging only because it
concentrates in an over-represented industry. With a single low-cardinality
categorical confounder, exact stratification is not an approximation of propensity
matching but the exact version of what it approximates, so the effect is directly
standardized,
\begin{equation}
  \bar e^{\,\text{std}} \;=\; \frac{1}{|G|}\sum_{g\in G} \bar e_g ,
\end{equation}
and the gap $\big||\bar e|-|\bar e^{\,\text{std}}|\big|/|\bar e|$ is reported as the
share of the apparent effect attributable to industry composition.

Aggregation is \emph{worst case} rather than average --- a desk choosing a model needs
the floor --- and maps to an ordinal rating
\begin{equation}
  \text{rating} =
  \begin{cases}
    1\ (\text{robust})   & \max_P \delta_P < 0.25,\\
    2\ (\text{moderate}) & 0.25 \le \max_P \delta_P < 0.60,\\
    3\ (\text{fragile})  & \text{otherwise.}
  \end{cases}
\end{equation}

\begin{keybox}[Why a rating and not another number]
The rating separates models that a single skill statistic ranks identically --- or
backwards. On a controlled panel an over-parameterized fit posts a \emph{higher}
in-sample-flavoured rank-IC ($0.78$ vs $0.34$) than a regularized one while rating
$3$ against $1$: it memorizes, and small realistic data defects destroy it.
\end{keybox}

\subsection{Sub-population skill}
The headline is an average over a heterogeneous universe and can hide a sign flip.
Skill is therefore reported per bucket under two cuts that ask different questions:
GICS sector and industry group (\emph{where} does this work, in the terms a PM
allocates by) and within-date quintiles of the point-in-time Fama--French loadings
$\beta_{\text{mkt}},\beta_{\text{smb}},\beta_{\text{hml}},\beta_{\text{mom}}$
(\emph{what kind of company} does it work on, in the terms a risk model prices). Each
bucket reports $\overline{\RankIC}$, its $t$-statistic, $R^2_{\text{oos}}$ and the
decile spread; buckets below a minimum name count are computed but flagged, because a
headline resting on one is a finding rather than a result.

On the live US panel the two cuts are informative in exactly this way: an incumbent
specification whose overall rank-IC is $-0.013$ carries $+0.040$ in Industrials, and
its skill concentrates in the lowest market-beta quintile.

\subsection{Champion selection and combination}
Candidates surviving validation compete against an equal-weight combination of
themselves, formed by averaging \emph{within-date} percentile ranks,
\begin{equation}
  \hat\alpha^{\text{ens}}_{i,t}
  \;=\; \frac1M\sum_{m=1}^{M}\frac{\rank_i\big(\hat\alpha^{(m)}_{i,t}\big)}{N_t},
\end{equation}
which is scale-free and so combines a GBDT and an elastic net on equal terms. The
persisted ensemble instead sums member predictions standardized by their own training
moments, because the serving path scores a bare array with no date grouping available;
the two agree on ordering closely enough that selection remains meaningful.

\section{Rolling Fama--French factor regression}
\label{sec:ffreg}
\emph{Code:} \code{xbrl\_sec/sec/sources/factor\_model.py}. Per ticker, over a
rolling $252$-day window (stepped every $21$ trading days, minimum $126$
observations), excess returns are regressed on the FF factors of the chosen model
(FF3/FF4/FF5/FF6):
\begin{equation}
  r_{i,t}-r_{f,t}=\alpha_i+\sum_k\beta_{i,k}\,f_{k,t}+\varepsilon_{i,t}.
\end{equation}
Factor series are region-resolved (US, JP, Developed, Emerging) from the available
Ken French datasets and converted to decimals; the dependent window is winsorized
to the $1$st/$99$th percentile.

\subsection{Robust estimation (Huber IRLS)}
Rather than plain OLS, coefficients are estimated by iteratively reweighted least
squares with Huber weights, which down-weight outlier days. Over a small number of
iterations, with scale $\hat s=\operatorname{median}(|e|)/0.6745$ and
$u=|e|/\hat s$,
\begin{equation}
  \hat\beta^{(k+1)}=\big(X\tr W^{(k)}X\big)^{-1}X\tr W^{(k)}y,\qquad
  w_i=\begin{cases}1,& u_i\le\delta,\\ \delta/u_i,& u_i>\delta,\end{cases}
  \quad \delta=1.345 .
\end{equation}

\subsection{HAC (Newey--West) standard errors}
Standard errors are heteroskedasticity- and autocorrelation-consistent with
Bartlett weights over $L=5$ lags:
\begin{equation}
  \widehat{\Var}(\hat\beta)=(X\tr X)^{-1}\hat S\,(X\tr X)^{-1},\qquad
  \hat S=\hat\Gamma_0+\sum_{\ell=1}^{L}w_\ell\big(\hat\Gamma_\ell+\hat\Gamma_\ell\tr\big),\quad
  w_\ell=1-\frac{\ell}{L+1},
\end{equation}
with $\hat\Gamma_\ell=\sum_t (x_t\varepsilon_t)(x_{t-\ell}\varepsilon_{t-\ell})\tr$.

\subsection{Stored diagnostics and derived series}
Each window stores $\alpha$ (annualized $\times252$), the betas, $t$-statistics,
$p$-values, and standard errors, plus $R^2$, adjusted
$R^2=1-(1-R^2)\tfrac{n-1}{n-k}$, $F$-statistic and its $p$-value, RMSE, residual
volatility $\hat\sigma_\varepsilon\sqrt{252}$ (the idiosyncratic $D$ of
\S\ref{sec:ffcov}), the Durbin--Watson statistic
$\sum_t(\varepsilon_t-\varepsilon_{t-1})^2/\sum_t\varepsilon_t^2$, and the design
condition number. A composite \code{quality\_score}$\in\{0,\dots,5\}$ counts how
many of five gates a window passes (adj.\ $R^2>0.20$; $F$ $p<0.05$; $n\ge200$;
condition number $<30$; $1.5\le \mathrm{DW}\le2.5$). Between refits, per-day
model-implied returns $\hat r_{i,t}=\alpha_i/252+\sum_k\beta_{i,k}f_{k,t}+r_{f,t}$
are written to \code{fact\_factor\_implied\_returns}. These loadings feed the WACC
(\S\ref{sec:wacc}) and the FF structured covariance (\S\ref{sec:ffcov}).

\section{Macro-cycle regime models}
\label{sec:regime}
\emph{Code:} \code{xbrl\_sec/sec/cycle/}. The regime stack has three layers: a
multimodal feature set, a dimensionality-reduction \emph{baseline} that produces
latent cycle factors, and two \emph{regime models} (HMM and VAE) that run on top,
each mapped to interpretable business-cycle phases and then linked to forecasting
through regime-conditioned factor IC.

\subsection{Feature construction}
\code{cycle/features.py} assembles a monthly, point-in-time, multimodal feature
matrix across three modalities --- \emph{macro} (growth, inflation, rates,
liquidity, credit), \emph{market}, and \emph{fundamental} --- each observation
stamped with its true availability date. Features are standardized with an
\emph{expanding} (walk-forward) mean and standard deviation,
$z_{f,t}=(x_{f,t}-\mu_{f,\le t})/\sigma_{f,\le t}$ (minimum $12$ observations), so
no future information leaks into any month's normalization.

\subsection{Baseline: PCA / dynamic-factor compression}
\code{cycle/baselines.py} compresses the (mean-centered) feature matrix to $K$
latent cycle factors via the SVD $X=U\Sigma V\tr$, taking the leading components
$\text{factors}=U_{:,1:K}\,\Sigma_{1:K}$ with explained-variance ratios
$\sigma_k^2/\sum\sigma^2$. The \code{dfm} (dynamic-factor) variant additionally
smooths the factors with an exponentially weighted moving average (span $3$),
approximating the persistence of a dynamic factor model. The factors are written
to \code{fact\_cycle\_state\_monthly.latent\_cycle} and are the input to the two
regime models.

\subsection{Phase calibration and labeling}
\label{sec:phasecal}
Latent factors have arbitrary sign and order, so \code{cycle/phase.py} anchors
them to interpretable phases. It (i) selects the factor most correlated with an
external stress/recession proxy from the macro store and orients it so ``high $=$
stress''; (ii) maps the factor to a rolling stress \emph{percentile} $s_t$ over a
$120$-month window; and (iii) thresholds $s_t$ into four phases
$\{$contraction, late\_cycle, mid\_expansion, early\_expansion$\}$. For US models
the thresholds are calibrated against the NBER recession chronology (via
\code{USREC}'s trough convention) by maximizing balanced accuracy
$\tfrac12(\mathrm{TPR}+\mathrm{TNR})$ over a grid; labels are then smoothed
(EWMA) and passed through a minimum-phase-duration filter that dissolves runs
shorter than three months. Phase confidence is a bounded function of distance from
the $0.5$ percentile,
$c_t=\clip(1.5\,|s_t-0.5|+0.45,\,0.45,\,0.90)$.

\subsection{Gaussian Hidden Markov Model}
\label{sec:hmm}
\code{cycle/hmm.py} fits a Gaussian HMM (default $K=4$ states, full covariance)
over the latent cycle factors $\xv_t$:
\begin{equation}
  \xv_t\mid s_t{=}k\sim\mathcal N(\boldsymbol\mu_k,\Sigma_k),\qquad
  A_{ij}=\Prob(s_t{=}j\mid s_{t-1}{=}i),\qquad
  \pi_k=\Prob(s_1{=}k).
\end{equation}
The parameters $\{\pi,A,\boldsymbol\mu_k,\Sigma_k\}$ are fit by Baum--Welch
(the EM algorithm for HMMs). The forward--backward recursions define
$a_t(k)=\Prob(\xv_{1:t},s_t{=}k)$ and $b_t(k)=\Prob(\xv_{t+1:T}\mid s_t{=}k)$,
from which the smoothed state posterior is
\begin{equation}
  \gamma_t(k)=\Prob(s_t{=}k\mid\xv_{1:T})
  =\frac{a_t(k)\,b_t(k)}{\sum_{k'}a_t(k')\,b_t(k')} .
\end{equation}
The per-month confidence is $\max_k\gamma_t(k)$ and the MAP state path is recovered
by Viterbi. States are mapped to phases by ordering their mean stress score
(\S\ref{sec:phasecal}). If \code{hmmlearn} is unavailable, a deterministic
$k$-means-style assignment with empirical transition counts produces the same
output shape.

\subsection{Temporal multimodal variational autoencoder}
\label{sec:vae}
\code{cycle/vae.py} is an optional VAE with encoder
$q_\phi(\zv\mid\xv)=\mathcal N\!\big(\boldsymbol\mu_\phi(\xv),\diag\boldsymbol\sigma_\phi^2(\xv)\big)$,
decoder $p_\theta(\xv\mid\zv)$, and prior $p(\zv)=\mathcal N(\mathbf0,I)$. It is
trained by maximizing the evidence lower bound with a temporal smoothness penalty:
\begin{equation}
  \mathcal L(\theta,\phi)=
  \underbrace{\E_{q_\phi}\!\big[\log p_\theta(\xv\mid\zv)\big]}_{\text{reconstruction}}
  -\;\beta\,\underbrace{\mathrm{KL}\!\big(q_\phi(\zv\mid\xv)\,\Vert\,p(\zv)\big)}_{\text{regularization}}
  -\;\eta\sum_t\lVert\boldsymbol\mu_{\phi,t}-\boldsymbol\mu_{\phi,t-1}\rVert^2,
\end{equation}
using the reparameterization
$\zv=\boldsymbol\mu_\phi+\boldsymbol\sigma_\phi\odot\boldsymbol\epsilon$,
$\boldsymbol\epsilon\sim\mathcal N(\mathbf0,I)$ (in the shipped configuration
$\beta=0.05$, $\eta=0.02$). The latent dimensions are split into interpretable
\emph{growth / inflation / rates-liquidity / credit-stress / market /
fundamentals} factors. The smoothed reconstruction-error stress percentile is
thresholded into phases exactly as in \S\ref{sec:phasecal}, and per-modality
reconstruction MSE attributes stress to macro vs.\ market vs.\ fundamentals. If
PyTorch is absent, a PCA surrogate produces the same persisted output shape.

\subsection{Regime-conditioned factor IC}
\label{sec:regic}
\code{cycle/ic.py} is the link from the regime model to forecasting. For each
fundamental factor $m$, \emph{conditioned on the regime}, it computes the rank
information coefficient against forward returns over each month and horizon:
\begin{equation}
  \IC_{\text{regime}}(m)=\rho_{\text{Spearman}}\big(x_{m,i},\,r^{\text{fwd}}_i \;\big|\; \text{regime}\big).
\end{equation}
Two variants are stored to \code{fact\_equity\_factor\_ic\_regime}: a
\emph{hard-assigned} IC (grouping by the MAP phase label) and a
\emph{probability-weighted} IC that weights each observation by the regime
posterior $\gamma_t(k)$ (soft assignment), using the weighted Spearman correlation
of the value and forward-return ranks. The result answers \emph{which factors have
actually predicted returns in conditions like today's}, and drives the scanner's
factor-family weights (\S\ref{sec:interest}).

\begin{modelbox}{From regime to score: the one live consumption path}
\small
\code{ic\_weights.family\_weights} reads the \emph{latest date's} regime-IC rows,
aggregates per-metric $|\IC|$ to family-level emphasis
$e_{\text{fam}}=\overline{|\IC|}_{\text{fam}}$ (normalized to sum to one over
value / growth / quality), and caches it. The scanner then re-splits its
fundamental score budget in proportion to $e_{\text{fam}}$: the interest score
leans harder on whichever factor family the current regime says is most
predictive, and reverts \emph{exactly} to the fixed weights when the IC table is
unavailable. The learned alpha model itself is regime-\emph{agnostic} --- it
learns from factor values, and the regime enters forecasting through the score
weights, not (currently) as a model feature.
\end{modelbox}

%======================================================================
\part*{Part IV\quad Deterministic valuation and scoring}
\addcontentsline{toc}{section}{Part IV\quad Deterministic valuation and scoring}

\section{Valuation models}
\label{sec:val}
\emph{Code:} \code{ai\_analyst/dcf\_engine.py}, \code{committee/wacc.py},
\code{committee/valuation.py}, \code{committee/segments.py}. Deterministic, on a
$7$-year explicit horizon (\code{DCF\_HORIZON\_YEARS}). The models are intentionally
transparent so the equity-research ``why'' of every number reads off the projected
statements; the LLM committee supplies only the bounded assumptions, and the math
below turns them into figures.

\subsection{Discounted cash flow (unlevered FCFF)}
For each forecast year $i=1,\dots,7$:
\begin{align}
  \text{Rev}_i&=\text{Rev}_{i-1}(1+g_i), &
  \text{EBIT}_i&=\text{Rev}_i\cdot m, \\
  \text{NOPAT}_i&=\text{EBIT}_i-\max(\text{EBIT}_i,0)\,\tau, &
  \text{FCF}_i&=\text{NOPAT}_i+\text{D\&A}_i-\text{Capex}_i-\Delta\text{NWC}_i,
\end{align}
with D\&A, Capex, and $\Delta$NWC taken as fractions of revenue (D\&A\% from the
latest historical ratio). Terminal value uses Gordon growth, and the
enterprise/equity bridge is
\begin{equation}
  \text{TV}=\frac{\text{FCF}_7(1+g_\infty)}{\text{WACC}-g_\infty},\qquad
  \text{EV}=\sum_{i=1}^{7}\frac{\text{FCF}_i}{(1+\text{WACC})^i}+\frac{\text{TV}}{(1+\text{WACC})^7},
\end{equation}
$\text{Equity}=\text{EV}-\text{NetDebt}$, per share $=\text{Equity}/\text{shares}$.
A $5\times5$ sensitivity grid varies WACC by $\pm200$\,bp and $g_\infty$ by
$\pm100$\,bp (a growth$\times$WACC grid is also produced around the base case).

\subsection{Cost of capital (WACC)}
\label{sec:wacc}
Cost of equity is built from the FF factor betas (\S\ref{sec:ffreg}) and long-run
premia $\Pi_k=252\cdot\overline{f_k}$:
\begin{equation}
  R_e^{\text{CAPM}}=r_f+\beta_{\text{Mkt}}\,\text{ERP},\qquad
  R_e^{\text{FF5}}=r_f+\sum_k\beta_k\,\Pi_k,
\end{equation}
with $r_f$ the 10Y Treasury (\code{FRED:DGS10}); both are clamped to
$[r_f+2\%,\,r_f+12\%]$. Cost of debt is $R_d=r_f+\text{spread}$, where the spread
comes from a Damodaran-style synthetic-rating table keyed on interest coverage
$\text{EBIT}/\text{Interest}$, blended with the realized $\text{interest}/\text{debt}$
rate when that looks plausible. Then
\begin{equation}
  \text{WACC}=\frac{E}{V}R_e+\frac{D}{V}R_d(1-\tau),\qquad V=E+D,
\end{equation}
$E$ market cap, $D$ total financial debt, clamped to $[6\%,16\%]$. The headline
uses $R_e^{\text{CAPM}}$ (the market convention and a stable discount rate); the
multi-factor $R_e^{\text{FF5}}$ is reported alongside as an exhibit. For the
sum-of-the-parts DCF, \code{segment\_wacc} shifts the equity beta and adds a
growth-risk premium so higher-growth segments carry a higher discount rate.

\subsection{Reverse DCF, scenarios, sum-of-the-parts, triangulation}
\emph{Reverse DCF} solves, by bisection, the flat growth $g^\star$ (or steady
margin $m^\star$) such that $\text{PerShare}(g^\star)=\text{price}$ --- i.e.\ what
the market is currently pricing in. \emph{Scenario weighting} runs the DCF once per
upside/base/downside scenario (sharing one deterministic share count and net debt)
and combines the per-share values $V_s$ with probabilities $\sum_s p_s=1$ into
$\text{FV}=\sum_s p_s V_s$. \emph{Sum-of-the-parts} values each reportable segment
on its own growth/margin and a segment-specific WACC, with segment FCF
$\approx\text{NOPAT}\times\text{conversion}$, and sums the segment enterprise
values. \emph{Triangulation} builds a football field from the SOTP range
(declared primary), the consolidated DCF scenario range, and peer-multiples-implied
values, reporting implied upside $\text{FV}/\text{price}-1$.

\section{Scoring and ranking models}
\label{sec:score}

\subsection{Value $+$ sentiment interest score}
\label{sec:interest}
\emph{Code:} \code{api/routers/screener\_agent.py}. Each component is min--max
normalized across the shortlist to $[0,1]$ ($v_{\text{PE}}$ inverted --- cheaper is
better), plus a normalized alpha $v_\alpha$ from the qlib model:
\begin{equation}
  \text{score}=w_\alpha\,v_\alpha+(1-w_\alpha)\sum_k w_k\,v_k,
  \qquad \text{interest}=100\cdot\text{score},
\end{equation}
with $w_\alpha\approx0.40$ (dominant when the alpha model covers the shortlist).
The fundamental sub-weights $w_k$ (FCF-yield, P/E, growth, MD\&A tone, news) start
from availability-branched defaults and are re-split between the \emph{value} and
\emph{growth} blocks by the current regime's factor IC (\S\ref{sec:regic}): with
family emphases $e_{\text{val}},e_{\text{grw}}$, the fundamental budget is allocated
$\propto e_{\text{val}}:e_{\text{grw}}$. With no model or IC available it reduces
exactly to the fixed-weight composite.

\subsection{Group relative-value composite}
\label{sec:group}
\emph{Code:} \code{ai\_analyst/committee/group.py}. For each metric $k$ (P/E,
EV/EBITDA, P/B, FCF yield, dividend, revenue growth/CAGR, gross/operating margin),
a cross-sectional $z$-score $z_{k,i}=(x_{k,i}-\bar x_k)/\hat\sigma_k$ and a signed
weight $w_k$ (negative for cheaper-is-better multiples). Loss-making multiples
(negative P/E, EV/EBITDA, P/B) are excluded from their metric's population rather
than ranked as ``cheap''. The composite is
\begin{equation}
  \text{composite}_i=\sum_k w_k\,z_{k,i},
\end{equation}
ordered descending; terciles map to attractive / fair / expensive. The per-name
audit trail $(x\to z\to \times w\to\text{contribution})$ sums exactly to the
composite. An optional LLM narrative may be attached; the ranking itself is
deterministic.

\subsection{MD\&A management tone}
An LLM reads an excerpt of each name's latest MD\&A and returns a structured tone
in $[-1,1]$, a positive/neutral/negative guidance label, and risk flags; the tone
feeds $v_{\text{tone}}$ in \S\ref{sec:interest}.

%======================================================================
\part*{Part V\quad How the models connect}
\addcontentsline{toc}{section}{Part V\quad How the models connect}

\section{The model map}
The three quantitative paths and the valuation path meet at the investment
committee. The valuation path answers what a business is worth; the quant path
answers what to expect and how to size it; the macro-cycle path conditions the
scoring on the regime; and the committee fuses the deterministic evidence and the
quant signals into a reasoned, auditable verdict.

\begin{figure}[t]
\centering
\begin{tikzpicture}[node distance=9mm and 14mm, every node/.style={font=\sffamily},
   mm/.style={proc, minimum width=34mm}]
  \node[store, minimum width=56mm] (wh) {Warehouse\;\;\scriptsize prices · fundamentals · factors · macro};

  \node[mm, below=11mm of wh, xshift=-52mm] (ff) {FF rolling regression\\\scriptsize $\beta$, residual vol};
  \node[mm, below=11mm of wh] (cyc) {Macro-cycle model\\\scriptsize HMM / VAE $\to$ regime};
  \node[mm, below=11mm of wh, xshift=52mm] (panel) {PIT fundamentals panel\\\scriptsize $z$-scored features};

  \node[mm, below=9mm of ff] (wacc) {WACC / structured $\Sigma$};
  \node[mm, below=9mm of cyc] (regic) {Regime-conditioned IC};
  \node[llm, below=9mm of panel, minimum width=34mm] (alpha) {qlib LightGBM alpha $\to\ \muv$};

  \node[mm, below=9mm of wacc] (dcf) {DCF · SOTP · reverse};
  \node[mm, below=9mm of regic] (interest) {Interest score / family wts};
  \node[llm, below=9mm of alpha, minimum width=34mm] (sigma) {Structured cov $\to\ \Sigma$};

  \node[llm, below=9mm of sigma, minimum width=34mm] (opt) {Portfolio optimizer $\to\ \wv$};

  \node[term, below=18mm of interest, minimum width=92mm] (comm)
    {Investment Committee \;\; \scriptsize evidence $+$ quant signals $\to$ memo $+$ fair value};

  \draw[arr] (wh) -- (ff); \draw[arr] (wh) -- (cyc); \draw[arr] (wh) -- (panel);
  \draw[arr] (ff) -- (wacc); \draw[arr] (cyc) -- (regic); \draw[arr] (panel) -- (alpha);
  \draw[arr] (wacc) -- (dcf); \draw[arr] (regic) -- (interest); \draw[arr] (alpha) -- (sigma);
  \draw[arr] (alpha.west) to[out=180,in=60] (interest.north east);
  \draw[arr] (sigma) -- (opt);
  \draw[arr] (dcf.south) |- (comm.west);
  \draw[arr] (interest.south) -- (comm.north);
  \draw[arr] (opt.south) |- (comm.east);
\end{tikzpicture}
\caption{The model map. Statistical, macro, and valuation paths converge on the
committee, which produces a reasoned valuation with its evidence rather than a
single label.}
\end{figure}

\vfill
\begin{center}\small\color{muted}
\emph{Independent, systematic research for educational purposes. Estimates are
model outputs, not investment advice.}
\end{center}

\end{document}
